\Needspace{15\baselineskip}
\section{The Liturgy of the Word}

The Collect's Amen does not introduce a religious lecture placed beside the sacrament. It opens the first of the Mass's two principal liturgies. In the readings God speaks to his people; Christ, present in his Word, proclaims the Gospel; the Homily opens that Word; silence and Psalm let it enter the heart; Creed and intercession answer it (GIRM 29, 55--71). The Liturgy of the Word and the Liturgy of the Eucharist form one act because the Christ who addresses the Church gives himself sacramentally within the same covenantal assembly (GIRM 28; \emph{Sacrosanctum Concilium} 7, 24, 48, 51, 56).

The distinction of modes remains essential. Christ's real presence in proclamation is not a lesser metaphor, but neither is it the substantial and continuous presence under the Eucharistic species. The Church venerates Scripture as she receives the one bread of life from the table of God's Word and Christ's Body (\emph{Dei Verbum} 21), without identifying book, voice, consecrated species, and incarnate Word. The ritual differentiations---ambo and altar, reader and Deacon, sitting and standing, proclamation and consecration---protect the unity from becoming confusion.

Silence belongs to this liturgy's structure. Before it begins, after readings, and after the Homily, measured pauses allow the Holy Spirit to join hearing with inward response (GIRM 45, 56). Such silence is no gap to be filled by another explanation: it lets a word not authored by the assembly judge, console, and convert.

Posture gives this sequence a bodily syntax. The faithful sit for the readings before the Gospel, the Responsorial Psalm, and the Homily; they stand for the Gospel acclamation and proclamation, the Creed, and the Universal Prayer (GIRM 43). Sitting is attentive receptivity rather than passivity, and standing is readiness rather than display. The changes let the Gospel's ceremonial summit and the Church's professing and interceding answer become visible.

\ordermovement{\latin{Lectio prima} (First Reading)}{A Received Pericope Becomes Present Address}{Proper biblical slot; ordinarily from the Old Testament on Sundays outside Easter Time, with the arrangement varying by celebration}{GIRM 29, 55--60, 128, 357; \latin{Ordo Lectionum Missae}}

A reader goes to the ambo and proclaims the appointed passage from the Lectionary; all sit and listen. At the conclusion the reader gives the received acclamation, and the people answer. The function is ministerial rather than presidential: a reader normally proclaims the first reading, while the Priest takes it only when no suitable reader is present (GIRM 58--59, 128). Book, ambo, reader, assembly, and answer make proclamation a public ecclesial action. Quietly following the text may assist hearing, but it is not the ritual act that the printed book assigns to the congregation.

The Lectionary's incipit may supply a name or circumstance omitted from the excerpt, yet the pericope remains a window cut from a complete biblical work. Historical prophet, narrator, sage, apostle, and evangelist do not become one generic ``liturgical voice.'' Serious mystagogy must therefore return the appointed passage to its literary argument, original hearers, canonical setting, and relation to Christ. GIRM 57 forbids replacing a reading with a non-biblical text precisely because the Church is receiving revelation, not choosing an inspirational theme.

Nehemiah gives a durable grammar for such reception. The people assemble, the book is opened, ears attend, interpretation is given, grief is redirected toward holy joy, and those without provision are remembered (Neh.~8:1--12). In Nazareth Jesus stands to read Isaiah, sits to interpret it, declares fulfillment, and exposes the resistance of familiar hearers (Luke~4:16--30). Hebrews describes God's Word as living, active, and discerning the heart rather than as inert material awaiting the listener's approval (Heb.~4:12--13). Proclamation makes no pericope newly inspired; it places the inspired Word in the Church's present act of obedient hearing.

Justin Martyr's mid-second-century Sunday outline already joins gathering, extended readings from the ``memoirs of the apostles'' and the prophets, and the president's exhortation (\emph{First Apology} 67). He does not describe the present three-year Lectionary, a fixed first-reading slot, or the modern ambo. His witness is more basic: biblical reading belongs to the Eucharistic assembly's inherited structure. Later Roman \emph{comites} and lectionaries organized pericopes according to the year; the medieval full Missal often concentrated reading in the celebrant. The current order differentiates the reader again and realizes Vatican II's command that the treasures of Scripture be opened more lavishly (\emph{Sacrosanctum Concilium} 51).

On Sundays and solemnities the first reading normally comes from the Old Testament, while during Easter Time Acts supplies the first witness; other days and ritual settings have their own arrangement (GIRM 357). This architecture refuses a Christianity detached from Israel. Promise, covenant, creation, wisdom, judgment, exile, and return are not a discarded prologue. They are the history in which the Paschal mystery is intelligible. Yet typology must respect the literal voice it fulfills. The Church hears Isaiah as Isaiah before making every prophetic phrase a direct prediction, and she hears Israel's lament as Israel's lament even when Christ and his Body pray it in a fuller key.

The people's concluding answer receives the reading as gift before every obscurity has been resolved. It does not canonize the reader's delivery, an introductory sentence, or a later interpretation. It thanks God for the Word entrusted through a human minister. The voice passes; the Word remains. Augustine's distinction between John the Baptist as a temporal voice and Christ as the abiding Word offers the right analogy: the voice serves precisely by sounding what it did not originate (\emph{Sermon} 293.3).

\ordermovement{\latin{Psalmus responsorius} (Responsorial Psalm)}{The Word Heard Becomes the Church's Answer}{Proper biblical chant integral to the Liturgy of the Word; limited approved alternatives; non-biblical substitution forbidden}{GIRM 37a, 61, 129; \latin{Ordo Lectionum Missae}}

The Psalm is not music between readings. GIRM 61 calls it an integral part of the Liturgy of the Word with great liturgical and pastoral importance because it fosters meditation upon the Word. The psalmist sings the verses from the ambo or another suitable place; the assembly ordinarily remains seated, listens, and returns the response. Singing at least the response is preferred. If singing is impossible, recitation must still serve meditation rather than haste.

This responsorial form creates a distinctive agency. The assembly does not answer revelation first with its own analysis. It answers one biblical passage with inspired poetry from another. The repeated response enters memory while the verses enlarge, correct, or complicate it. A cry of abandonment, royal promise, wisdom warning, penitential plea, or hymn of creation cannot be reduced to mood music. The actual Psalm and the first reading determine whether the relation is verbal echo, theological fulfillment, contrast, lament, thanksgiving, or sustained petition.

The Psalter itself teaches this conversion of hearing into speech. Psalm~95 moves from jubilant summons to kneeling worship and then lets God's warning break into the song: today, if his voice is heard, the heart must not harden. Colossians commands the Word of Christ to dwell richly through psalms, hymns, and spiritual songs (Col.~3:16). The response is therefore contemplative and ethical at once. Lips that repeat a verse are being schooled in a judgment they may not yet feel and a hope they did not invent.

Athanasius observes that the Psalms differ from other biblical books because the singer appropriates their words as his or her own movements of soul; they become a mirror and a school of fitting utterance (\emph{Letter to Marcellinus} 10--12). Augustine's reading of the \emph{totus Christus} adds a Christological and ecclesial depth: in the Psalms Head and members can speak within one voice without erasing the historical prayer of Israel (\emph{Enarrationes in Psalmos} 85.1). These patristic principles explain liturgical appropriation, but they do not pre-decide what a particular Psalm means. Exegesis begins with its own poem.

Responsorial psalmody is ancient, while its concrete Roman history is layered. Early Christian and late antique witnesses attest psalm singing with congregational responses; the Roman Gradual eventually concentrated the post-reading chant in a highly elaborated form. The postconciliar Lectionary restored a responsorial Psalm closely coordinated with each reading without forbidding the received Gradual books. In the United States GIRM 61 admits the Responsorial Gradual from the \emph{Graduale Romanum}, forms from the \emph{Graduale Simplex}, and approved antiphon-and-Psalm collections. It then sets a sharp boundary: songs or hymns may not replace the Responsorial Psalm. The restriction protects not a preferred musical style but the inspired answer's place inside the proclaimed Word.

\ordermovement{\latin{Lectio secunda} (Second Reading)}{The Apostolic Witness in Its Own Voice}{Proper biblical slot where appointed, notably on Sundays and solemnities}{GIRM 57--59, 130, 357; \latin{Ordo Lectionum Missae}, Introduction 66--68}

Where three readings are appointed, a reader proclaims the second from the ambo and the people receive it with the same concluding acclamation. On Sundays and solemnities it normally comes from an Apostle---a New Testament letter or, according to the season, Revelation---between the Old Testament reading and the Gospel (GIRM 357; \emph{Ordo Lectionum Missae}, Introduction 66). ``Second'' denotes ritual order, not secondary authority.

The apostolic writings were composed for concrete churches, coworkers, crises, and hopes. Paul commands a letter sent to Colossae to be read also at Laodicea and another to be exchanged in return (Col.~4:16); he solemnly charges that First Thessalonians be read to all the brothers and sisters (1~Thess.~5:27). A local address thereby becomes trans-local ecclesial instruction. Liturgical proclamation continues that passage without pretending that Corinth, Rome, or the seven churches of Revelation were generic audiences. Their questions belong to the Word now received.

Justin's phrase ``memoirs of the apostles or writings of the prophets'' shows apostolic and prophetic witnesses already sounding together on Sunday (\emph{First Apology} 67). Near the end of the second century, Irenaeus argues that the apostolic proclamation is publicly preserved in the churches rather than available only through a secret elite (\emph{Adversus haereses} III.3.1--4). The second reading makes that public character audible. Its minister does not deliver private gnosis; the apostolic text addresses a Church standing within a visible communion.

The current Sunday Lectionary uses a three-year cycle and several principles of arrangement. In Ordinary Time an apostolic book is often read in a semi-continuous sequence, while the Old Testament reading is more directly chosen in relation to the Gospel (\emph{Ordo Lectionum Missae}, Introduction 66--68). Consequently a preacher or reader should not force an immediate thematic link that the Lectionary did not intend. The second reading may bring its own argument, rebuke, consolation, or ecclesiology into the celebration. Its place still joins the covenants: Israel's Scriptures, apostolic interpretation and mission, and the Gospel's climactic proclamation disclose one history centered upon Christ's Paschal mystery.

\ordermovement{\latin{Sequentia} and \latin{Acclamatio ante Evangelium}}{Poetry and Acclamation before the Lord Who Speaks}{Sequence conditional; Gospel acclamation seasonal and normally sung, with stated omission rules}{GIRM 37a, 62--64, 131; \latin{Ordo Lectionum Missae}}

Two distinct forms may stand here and must not be collapsed. A Sequence, where appointed, is sung before the Alleluia. It is obligatory on Easter Sunday and Pentecost and optional on the other days for which the Missal supplies one (GIRM 64). The Gospel acclamation then begins as an act in itself. All stand; choir or cantor leads, the people sing the acclamation, and the proper verse is given its ministerial voice (GIRM 62). Poetry unfolds a feast; acclamation welcomes the Lord about to speak.

Outside Lent the acclamation is Alleluia with its appointed verse. During Lent another verse or a Gradual tract takes its place. When there is only one reading before the Gospel, GIRM 63 gives several arrangements and permits the Alleluia or Lenten verse to be omitted if it is not sung. The provision should not be turned into a general license to recite a musical acclamation perfunctorily or to substitute a favorite refrain. The rite assigns the people's standing song the work of greeting Christ in his approaching Gospel.

Alleluia is Hebrew praise carried untranslated through Greek, Latin, and vernacular worship. Revelation~19:1--9 surrounds it with God's victory, the exposure of evil, the reign of the Lord, and the wedding feast of the Lamb. It is therefore neither a filler syllable nor a command to feel cheerful. Augustine calls it the Christian song of the heavenly homeland and insists that the pilgrim sings while still laboring on the road (\emph{Sermon} 256.1--3); elsewhere he gives its plain sense as praise of the Lord (\emph{Enarrationes in Psalmos} 148.1). The liturgical acclamation holds present pilgrimage and eschatological praise together as the Gospel book moves toward the ambo.

Gregory the Great's 598 letter to John of Syracuse shows that the extent of Roman Alleluia use was already a debated and differentiated custom; it does not make Gregory the inventor of the acclamation (\emph{Registrum epistolarum} IX.26). Medieval Western musical practice then generated the Sequence as a large poetic genre in relation to the Alleluia complex. Notker's late ninth-century preface to his \emph{Liber hymnorum} describes one route by which texts helped singers retain long melodic phrases, but his account is not a universal origin story for every Sequence. The genre developed far beyond mnemonic syllables into sustained biblical and doctrinal poetry.

The few Sequences retained in the Roman Missal demonstrate that each requires its own reading. Easter's poem stages Paschal witness; Pentecost's invokes the Spirit; Corpus Christi and the Sorrows of Mary have different histories, rhetoric, and conditions of use. None is itself a biblical reading, and no generic account of ``meditation'' can replace its text. Its position before Alleluia in the third typical edition clarifies the ascent: ecclesial poetry yields to the people's acclamation, and the acclamation opens toward the canonical Gospel.

\begin{studybox}{The Gospel acclamation is not commentary}
The proper verse may announce a theme or echo the Gospel, but the movement's primary direction is personal and liturgical: the assembly greets Christ who is about to address it. Its standing posture, repeated acclamation, and procession belong together.
\end{studybox}

\ordermovement{Preparation and Procession for the Gospel}{Purified Speech, Received Mission, and a Book in Motion}{Ministry-dependent blessing or quiet preparation; Book of the Gospels, candles, and incense used as provided}{GIRM 60, 131--133, 175}

If a Deacon proclaims the Gospel, he bows before the Priest and asks a blessing; the Priest invokes the Lord's aid over his heart and lips. The Deacon then takes the Book of the Gospels from the altar and processes to the ambo, preceded by incense and candles when they are used. If no Deacon is present, the Priest who will proclaim bows before the altar and says the quiet prayer beginning \latin{Munda cor meum}; he does not bless himself as though commissioning another minister. When the Book of the Gospels is on the altar, he then takes it and goes to the ambo with the ministers preceding him as provided (GIRM 131--133, 175). Preparation is interior, hierarchical, and spatial at once.

Isaiah~6:1--8 supplies the deep biblical pattern. Vision of divine holiness exposes unclean lips; fire from the altar cleanses; forgiveness precedes the prophet's answer to the question of sending. The liturgical prayer does not identify Deacon or Priest simply with Isaiah or the thurible with the prophet's coal. It receives the scriptural order: no herald treats the holy Word as personal possession, and purification serves mission. Romans makes the same ministerial point from another direction: preaching presupposes being sent (Rom.~10:14--17).

The difference between Deacon and Priest is ritually exact. The Gospel belongs by office to the Deacon when he is present (GIRM 59, 94). He receives a blessing because his proclamation is service within an ordered Church; the Priest proclaims only in the Deacon's absence and prepares quietly. The people do not join the private prayer or the ministerial blessing. Their part is the standing acclamation already begun. Differentiation here does not create unequal access to Scripture; it gives the climactic proclamation an ecclesial minister rather than a self-appointed speaker.

Book, lights, and incense intensify honor without depreciating the first and second readings. GIRM 60 explicitly calls the Gospel the high point of the Liturgy of the Word and names the special signs that distinguish it. Light readily bears the Johannine confession that life in the Word is the light of humanity (John~1:4--9); incense recalls prayer rising before God (Ps.~141:2; Rev.~8:3--4). These texts illuminate the signs but do not prescribe candle number, route, or thurible. Ceremonial particulars come from the Church's received books.

By the early medieval period \emph{Ordo Romanus I} already shows a differentiated Roman Gospel procession with Deacon, book, lights, incense, and blessing. The current movement inherits that public architecture while simplifying or expanding according to actual ministers. A Book of the Gospels carried at the Entrance has rested upon the altar; now it travels from altar to ambo. The route makes the unity of Word and sacrifice visible without confusing their places. The book is honored as the instrument of proclamation, then opened so that the written witness may become a living ecclesial address.

Ignatius calls Deacons ministers of the mysteries of Jesus Christ and servants of God's Church (\emph{Letter to the Trallians} 2.3). That early witness does not document the modern Gospel procession, but it keeps diaconal proclamation from being reduced to a convenient reading assignment. The Deacon's received blessing and public bearing of the Gospel express a ministry ordered to Christ's mystery within the Church.

\ordermovement{\latin{Evangelium} (Gospel Proclamation and Veneration)}{Christ Speaks through a Ministerial Voice}{Proper Gospel in a fixed dialogical and ceremonial frame; proclaimed by Deacon or Priest}{GIRM 29, 58--60, 134, 175}

At the ambo the Deacon or Priest greets the people, announces the Gospel, signs the book and then forehead, mouth, and breast; the people make the three small signs upon themselves. If incense is used, the book is incensed. All stand and listen. The proclamation ends in a dialogical acclamation; the minister kisses the book while saying the appointed quiet formula. When a Bishop presides, the book may be carried to him to be kissed, and he may bless the people with it according to the ceremonial (GIRM 134, 175).

The triple signing is a compact mystagogy of reception. The Gospel is to govern understanding, be confessed by the lips, and dwell in the inward person. The gesture does not mechanically protect three body parts or make private reading unnecessary. It submits thought, speech, and desire to the Word before a sentence has sounded. Standing similarly signifies readiness and honor; it does not imply that the other readings were less inspired while the people sat.

Luke's Nazareth scene shows written prophecy becoming present address in Christ's mouth (Luke~4:16--21). The risen Christ on the Emmaus road opens Moses and all the prophets through his own Pasch before he is recognized in the breaking of bread (Luke~24:25--35). His saying that reception of the apostolic messenger is reception of himself and the One who sent him gives ministerial proclamation its Christological seriousness (Luke~10:16). Yet the ordained reader is not another incarnation. Christ speaks through an ecclesially commissioned voice and canonical text; the minister remains accountable to both.

Irenaeus's famous account of the fourfold Gospel insists upon one Gospel under four forms and relates the evangelists to the one divine economy (\emph{Adversus haereses} III.11.8). Augustine likewise refuses to confuse the herald with the eternal Word: John is voice, while Christ is the Word who abides (\emph{Tractatus in Ioannem} 1.8--9). These distinctions make veneration intelligible. Incense, acclamation, and kiss are directed through book and proclamation toward Christ; they are not worship of paper or admiration of a performer's voice.

The ceremonial frame has developed. Justin attests apostolic memoirs being read on Sunday but gives no triple signing or book kiss (\emph{First Apology} 67). Late antique and early medieval churches increasingly distinguished the Gospel through its minister, procession, standing, lights, and incense; small signs and quiet prayers appear in later strata. The postconciliar rite retains these accumulated honors while restoring the ambo and public proclamation to their full function; where an approved vernacular is used, the hearing is correspondingly direct. It does not claim every gesture arose at once.

The concluding kiss has a penitential edge because its quiet prayer asks that the Gospel's words remove fault. The sign therefore cannot be reduced to congratulating the minister or closing a successful reading. The Word just acclaimed also judges. The assembly now sits not to leave that address behind but to receive its liturgical exposition.

\ordermovement{\latin{Homilia} and Silence}{The Proclaimed Mystery Opened toward Faith and Life}{Required at Sunday and holy-day Masses with a congregation absent grave reason; otherwise recommended as specified}{GIRM 29, 45, 55--56, 65--66, 136; CIC 767; \emph{Sacrosanctum Concilium} 35, 52}

The Homily is part of the liturgical action. It ordinarily belongs to the Priest celebrant, though he may entrust it to a concelebrating Priest or, when appropriate, to the Deacon; in particular cases a Bishop or Priest present but unable to concelebrate may preach. It is never assigned to a lay person at Mass (GIRM 66; CIC 767). That boundary does not deny the baptismal vocation to witness or teach. It locates this act of preaching within ordained presidency and the sacramental assembly.

Its matter is also received. The Homily explains some aspect of the readings or another text from the day's Ordinary or Proper, with attention to the mystery celebrated and the hearers' actual needs (GIRM 65). It is not licensed free association, partisan commentary, a detached theology lecture, or a second proclamation that displaces Scripture. The preacher serves a text and a mystery already given. Because those gifts address concrete lives, faithful preaching cannot be merely antiquarian exegesis either.

Jesus's synagogue discourse moves from Isaiah's written words to their fulfillment ``today'' and exposes the hearers' resistance (Luke~4:16--30). On the Emmaus road he interprets the Scriptures through the necessity of his suffering and glory, bringing despondent disciples toward recognition and mission (Luke~24:25--35). Justin describes the president, after the readings, admonishing and inviting the assembly to imitate the good things heard (\emph{First Apology} 67). The ancient homiletic act is thus neither information alone nor exhortation without exegesis: Word, Paschal fulfillment, hearer, and changed life meet.

Patristic preaching demonstrates immense variety rather than one ideal length or style. Origen follows difficult biblical sequences closely; Chrysostom joins textual argument to fierce moral application; Augustine moves between clause, doctrine, sacrament, and pastoral appeal. Gregory the Great formulates the enduring pastoral demand that one mode of exhortation does not suit hearers of different dispositions (\emph{Regula pastoralis} III, Prologue). Adaptation, however, concerns how the same faith reaches actual people; it does not authorize changing the Gospel to mirror every preference.

In the Roman tradition preaching did not simply vanish for a millennium, but its frequency, language, minister, and relation to the proper texts varied greatly. Vatican II required the Homily to be highly esteemed as part of the liturgy and forbade its omission on Sundays and holy days with a congregation except for serious reason (\emph{Sacrosanctum Concilium} 52). The current GIRM specifies a grave reason for omission and recommends preaching especially on weekdays of Advent, Lent, and Easter Time and when people gather in greater numbers (GIRM 66). This is a juridical and theological restoration of regular liturgical preaching, not the invention of Christian exposition.

Silence after the Homily completes rather than competes with it (GIRM 56, 66). The preacher relinquishes the last audible word so that the Spirit may press the proclaimed mystery into memory and will. Applause, immediate explanation, or hurried transition can make the Homily appear self-enclosed. The provided silence confesses that neither eloquence nor office causes the Word's fruit. Preaching serves; God gives the growth.

\ordermovement{\latin{Professio fidei} (Profession of Faith)}{The Many Speak the Church's One Faith}{Niceno-Constantinopolitan Creed when prescribed; Apostles' Creed admitted especially in Lent and Easter; baptismal promises replace it in appointed rites}{GIRM 67--68, 137; \latin{Ordo Missae}}

On Sundays and solemnities, and at other celebrations of a more solemn character when provided, Priest and people sing or say the Creed together (GIRM 68). The Niceno-Constantinopolitan Creed is the normal form. The Roman baptismal symbol called the Apostles' Creed may be used instead, especially during Lent and Easter Time; where a rite appoints renewal of baptismal promises, that renewal replaces the Creed. These alternatives are received and conditional. They are not raw material for an abbreviated or privately revised profession.

The movement answers the Word before the Eucharistic mystery begins (GIRM 67). Its first-person singular is corporate rather than individualistic: many mouths confess one faith into which each person was baptized. The Creed does not summarize only the day's readings. It traverses creation, the Father's sovereignty, the Son's Incarnation and Pasch, the Holy Spirit, the Church, Baptism, resurrection, and the world to come, locating each proper pericope within the whole economy of salvation.

Its scriptural density cannot be reduced to proof texts. Genesis~1 and John~1 belong to confession of creation; Luke~1 and John~1 to the Incarnation; 1~Corinthians~15 to Christ's resurrection and the promised resurrection of the dead; John~14--16 and Acts~2 to the Spirit; Ephesians~4:1--6 to the one Church and one Baptism; Revelation~21--22 to the awaited consummation. The Creed is the Church's rule for hearing these witnesses together, not a rival source standing above Scripture.

Cyril of Jerusalem tells catechumens that the faith was not composed according to human whim but that its most necessary points were gathered from the whole of Scripture, like mustard seed holding many branches (\emph{Catechetical Lecture} 5.12). His local fourth-century symbol is not textually identical with either form in the present Missal; his image explains their function. A concise ecclesial rule guards the canonical whole in the believer's memory. Tertullian's earlier ``rule of faith'' likewise shows Christian confession operating as a public boundary against secret innovation, though it is not the later conciliar text (\emph{De praescriptione haereticorum} 13).

The Nicene Council of 325 and the Constantinopolitan Council of 381 are decisive stages in the received Creed's doctrinal history; Chalcedon in 451 cites the expanded symbol as the faith of the 150 Fathers. The Western confession of the Spirit's procession with the \latin{Filioque} spread later and entered Roman use within a long doctrinal and liturgical history (CCC 246--248). None of these facts means the Creed occupied its current place at fourth-century Roman Masses. Spain prescribed liturgical recitation at the Third Council of Toledo in 589 (canon~2), in a position different from Rome's; Rome received the Creed at Mass in 1014 under Benedict~VIII at Emperor Henry~II's request. Ancient doctrine and later Roman placement are separate historical claims.

The Apostles' Creed is ancient and baptismal in lineage, but its familiar title does not prove that the Twelve composed the received text clause by clause. Early Roman baptismal interrogations and rules of faith developed into the later symbol; medieval accounts of each Apostle contributing an article belong to reception history rather than demonstrable authorship. Its current Lenten and Easter suitability comes from this baptismal character, not from lesser doctrinal weight.

All stand. At the words confessing the Incarnation, all make a profound bow; on the solemnities of the Annunciation and Nativity they genuflect (GIRM 137). The body does not add a decorative devotion to abstract propositions. It confesses that the eternal Son truly took flesh in history. The gesture is memorial, not a repeated Incarnation. From that bow the Creed moves through Cross and resurrection toward the life of the world to come, schooling present bodies in the hope of their own resurrection.

\ordermovement{\latin{Oratio universalis} (Universal Prayer)}{A Baptismal Priesthood Intercedes for the World}{Normally present with a congregation; categories stable, intentions locally composed under law; ministerial roles differentiated}{GIRM 69--71, 138; \emph{Sacrosanctum Concilium} 53}

Having heard and professed the Word, the people exercise their baptismal priesthood by interceding for the salvation of all (GIRM 69). The Priest directs the prayer from the chair, introducing it and concluding with an oration. A Deacon normally announces the intentions when present; otherwise a cantor, reader, or another lay member of the faithful may do so from the ambo or another suitable place. The people stand and answer each intention with a common invocation or pray in silence (GIRM 71, 138). The prayer is therefore presidential, ministerial, and congregational in distinguishable ways.

Its usual horizon has four parts: the needs of the Church; public authorities and the salvation of the whole world; those burdened by any difficulty; and the local community (GIRM 70). A Confirmation, Marriage, funeral, or other particular celebration may focus the sequence more closely. The categories are stable; the intentions are composed with wise freedom, sobriety, and few words. They should name who or what is entrusted to God rather than become miniature homilies, policy arguments, news summaries, or project-authored sample prayers.

First Timothy supplies almost the whole theological span: supplications and thanksgivings are to be made for everyone, including rulers, because God wills all to be saved and Christ is the one mediator (1~Tim.~2:1--6). The Church prays for Peter under persecution (Acts~12:5), and Paul commands intercession for all the saints and for his own evangelical speech (Eph.~6:18--20). Universal prayer is thus neither inward-looking congregational housekeeping nor the pretense that the Church already governs the world. She carries before God needs she cannot master.

Justin places common prayer after the readings and exhortation in his Sunday outline; elsewhere he describes the baptized praying for themselves, the newly illumined, and people everywhere (\emph{First Apology} 65, 67). First Clement's great prayer asks holiness for the elect, help for the afflicted, return for the wandering, and wisdom and peace for earthly rulers (\emph{1 Clement} 59.2--61.3). These witnesses do not supply the modern four categories word for word, but they reveal the early breadth: ecclesial prayer extends beyond those physically present, even toward authorities who may be hostile.

The broad intercessions receded from the ordinary Roman Mass during the early medieval development, while related petitions survived in special solemn forms and other rites. Vatican II explicitly ordered their restoration, especially on Sundays and holy days, so that prayer would be made for Church, civil authorities, those in need, all humanity, and the salvation of the world (\emph{Sacrosanctum Concilium} 53). The current order realizes that mandate without pretending that every local set of intentions is ancient text.

This movement concludes the Liturgy of the Word by turning hearing outward. The assembly has not listened only for private enlightenment. God's promises and judgments create responsibility for the Church, rulers, suffering persons, neighbors, enemies, and the whole world. The Priest's concluding oration gathers those petitions without erasing the people's voice. Then the assembly sits and created gifts are brought forward. Intercession has already widened the local gathering to the dimensions of the sacrifice that will be offered for the living and the dead in the communion of the whole Church.
